\documentclass[parskip=half]{scrartcl}
\usepackage[ngerman]{babel}
\usepackage[utf8]{inputenc}
\usepackage{graphicx}
\usepackage{amsmath}
\usepackage{amssymb} % Diese Zeile hinzufügen
\usepackage{float}
% \usepackage[autostyle=true]{csquotes}
\usepackage{tabularx}
\usepackage{booktabs} % Für schönere Linien (optional, aber empfohlen)

% Zitierstil und Literaturverzeichnis
\usepackage[style=apa, backend=biber]{biblatex}
\DeclareLanguageMapping{ngerman}{ngerman-apa}

\addbibresource{Literatur.bib}


\title{Klassische Markowitz Optimierung versus HRP}
\author{Niclas Treib}

\begin{document}
\thispagestyle{empty}
	\begin{center}
		\includegraphics[width=0.5\textwidth]{HS_Koblenz_Logo.jpg}
		\vspace{2cm} % Erzeugt 1 cm vertikalen Abstand zwischen Logo und Titel
		\\
		% 2. Der Titel (in Fett und etwas größer)
		{\LARGE \textbf{Klassische Markowitz Optimierung versus HRP}}
		
		\vspace{2cm}
		
		{\Large \textbf {Bachelorarbeit}} \\
		
		\vspace{1cm}
		
		\large{im Studiengang Wirtschaftsmathematik \\
		Hochschule Koblenz, RheinAhrCampus,			\\
		Remagen}
		
		\vspace{1cm}
		
		vorgelegt von \\
		{\textbf {Niclas Treib}} \\
		geb. am 21.09.1999 \ in Bad Neuenahr-Ahrweiler
		
		\vspace{1.5cm}
		
	\end{center}

\begin{center}
\begin{tabular}{lll}
	\textbf{Erstprüfer:} & & Prof. Dr. Anselm Hudde \\
	\textbf{Zweitprüfer:} & & Prof. Dr. Martina Brück \\
	\\ % & springt in die nächste Spalte der Tabelle
	\textbf{Durchgeführt bei:} & & Hochschule Koblenz
\end{tabular}
\end{center}

\vspace{2cm}

\begin{center}
	\large {30. Juni 2026}
\end{center}
\clearpage
\setcounter{page}{1}
\pagenumbering{Roman}

\tableofcontents
\clearpage

\listoffigures
\clearpage

\listoftables
\clearpage

\section*{Symbolverzeichnis}



\clearpage

% \pagenumbering{arabic} <- arabische Zahlen als Nummerierung
\setcounter{page}{1}
\pagenumbering{arabic}

	Hierarchical Risk Parity
	\section{Einleitung}
	Nach mehr als sieben Jahrzehnten ist nach wie vor Markowitz das vorherrschende Optimierungsverfahren in der Portfolioallokation. Es liefert zwar gute Ergebnisse, doch weist auch signifikante Schwächen auf. Das Verfahren reagiert empfindlich auf Schätzfehler in den Eingangsparametern und es neigt dazu wenige schwach volatile Aktien stark zu gewichten. Das sind zwei Eigenschaften, die zu nicht optimalen Ergebnissen führen können und unter gewissen Marktsituation nicht gewollt sind. \\ \\
	Der Fortschritt im Bereich des Machine Learnings ermöglicht Verfahren, die diese Schwächen nicht mehr aufweisen oder wo diese Schwächen zumindest weniger stark ausgeprägt sind. Marcos Lopez de Prado hat das HRP-Verfahren (Hierarchical Risk Parity) entwickelt, welches Baummodelle nutzt um die Empfindlichkeit gegenüber den Schätzfehlern zu verbessern. Das kann insbesondere bei der Bestimmung eines Minimum-Varianz-Portfolios die Ergebnisse verbessern. \\ \\
	Meine Arbeit beschäftigt sich daher mit dem Vergleich der zwei genannten Modelle und der Frage, ob HRP bessere Ergebnisse liefert als die klassische Markowitz-Optimierung. \\
	Hierfür betrachte ich ein Portfolio bestehend aus 30 Aktien über die letzten 20 Jahre. \\
	
	
	\clearpage
	\section{Theoretische Grundlagen}
	Um das Optimierungsverfahren nach Markowitz und das HRP-Modell zu untersuchen ist es zunächst notwendig, die theoretischen Grundlagen zu klären.
	\subsection{Rendite}
	Die Rendite beschreibt die relative Wertänderung eines Investments bezogen auf den investierten Kapitalwert am Anfang des Betrachtungszeitraums:
	
	\begin{equation}
		R_t = \frac{S_t - S_{t-1}}{S_{t-1}} \label{eq:rendite}, \qquad t \in \{1, \dots, n\}.
	\end{equation} 
	Basierend auf der Rendite können wir die weiteren statistischen Werte ermitteln.
	
	\subsection{Erwartungswert der Rendite}
	Der Erwartungswert der Rendite wird klassischer Weise durch das arithmetische Mittel der Beobachtungen geschätzt. Für die Schätzung ergibt sich damit:
	
	\begin{equation}
		\hat{\mu}_i = \frac{1}{T} \sum_{t=1}^{T} R_{i,t}. \label{eq:erwartungswert}
	\end{equation}
	Wobei $T$ die Anzahl der Beobachtungen ist. \\
	
	\subsection{Varianz und Standardabweichung}
	Die Varianz misst die quadratische Abweichung vom Erwartungswert. Die Varianz wird in dieser Arbeit als $\hat{\sigma}^2$ oder mit $\operatorname{Var}()$ bezeichnet. Sie ist gegeben als
	
	\begin{equation}
		\hat{\sigma}_i^2 = \frac{1}{T-1} \sum_{t=1}^{T} (R_{i,t} - \hat{\mu}_i)^2. \label{eq:varianz}
	\end{equation}
	Durch $T-1$ im Nenner ist der Schätzer $\hat{\sigma}_i^2$ für die Varianz erwartungstreu (Bessel-Korrektur).
	
	Die Standardabweichung bzw. Volatilität ergibt sich als $\hat{\sigma}_i$.
		
	\clearpage
	
	\subsection{Kovarianz}
	Zur Analyse des gemeinsamen Bewegungsverhaltens zweier Assets $i$ und $j$ benötigen wir die Kovarianz. Sie beschreibt die gemeinsame Bewegung durch:
	
	\begin{equation}
		\operatorname{Cov}(R_i,R_j) = \frac{1}{T-1} \sum_{t=1}^{T} (R_{i,t} - \hat{\mu}_i)(R_{j,t} - \hat{\mu}_j). \label{eq:kovarianz}
	\end{equation}
	\textbf{Bemerkung:}
	Die Kovarianz eines Assets mit sich selbst entspricht der Varianz, denn
	
	\begin{equation}
		\operatorname{Cov}(R_i,R_i) = \frac{1}{T-1} \sum_{t=1}^{T} (R_{i,t} - \hat{\mu}_i)(R_{i,t} - \hat{\mu}_i) = \frac{1}{T-1} \sum_{t=1}^{T} (R_{i,t} - \hat{\mu}_i)^2 = \hat{\sigma}_i^2.
	\end{equation}
	% Variante 1: Abgesetzt und zentriert (Empfohlen für echte Formeln)
	% \[ x + y = z \]
	% \[ x \cdot y \, dx \]
	
	% Variante 2: Im Text, aber mit manuellem Umbruch
	% $xy = dx$ \\
	% $a^2 + b^2 = c^2$
	\subsection{Korrelation}
	Die Kovarianz hängt von der Skalierung der Renditen ab und ist somit nicht leicht zu interpretieren. Weshalb man als leichter zu interpretierendes Maß die Korrelation betrachtet. Sie ist auf das Intervall $[-1,1]$ normiert und berechnet sich wie folgt:
	
	\begin{equation}
	\rho_{i,j} = \frac{\operatorname{Cov}(R_i,R_j)}{\hat{\sigma}_i \cdot \hat{\sigma}_j} , \quad \text{mit } \rho_{i,j} \in [-1,1] \quad \forall \ i,j. \label{eq:korrelation}
	\end{equation}
	Wobei $\rho_{ij} = 0$ bedeutet, dass zwei Assets unkorreliert sind. Und $\rho_{ij} = \pm 1$ den perfekt korrelierten bzw. perfekt unkorrelierten Fall beschreibt.
	
	\textbf{Bemerkung:} Die Korrelation eines Assets mit sich selbst ist immer 1, denn:
	\begin{equation}
		\rho_{i,i} = \frac{\operatorname{Cov}(R_i,R_i)}{\hat{\sigma}_i \cdot \hat{\sigma}_i} = \frac{\operatorname{Var}(R_i)}{\hat{\sigma_i}^2} = \frac{\hat{\sigma_i}^2}{\hat{\sigma_i}^2} = 1.
	\end{equation}
	Das würden man auch intuitiv vermuten, denn eine Aktie korreliert immer perfekt mit sich selbst.
	
	\subsection{Kovarianz- und Korrelationsmatrix}
	Um die Kovarianz bzw. die Korrelation mehrerer Assets simultan darzustellen nutzt man Matrizen. Sie werden als Kovarianz bzw. Korrelationsmatrix bezeichnet.
	Sie haben die Form:
	
	% m x n Matrix, aber m = n weil die Matrix ja quadratisch ist!
	\begin{equation}
		\Sigma = \begin{pmatrix}
			\operatorname{Var}(R_1) & \operatorname{Cov}(R_1,R_2) & \cdots & \operatorname{Cov}(R_1,R_n) 	\\
			\operatorname{Cov}(R_2,R_1) & \operatorname{Var}(R_2) & \cdots & \operatorname{Cov}(R_2,R_n) 	\\
			\vdots & \vdots &  & \vdots										 								\\
			\operatorname{Cov}(R_n,R_1) & \operatorname{Cov}(R_n,R_2) & \cdots & \operatorname{Var}(R_n)
		\end{pmatrix},										\label{eq:kovarianzmatrix}
	\end{equation}
	
	\begin{equation}
		R = \begin{pmatrix}
			1 & \rho_{1,2} & \cdots & \rho_{1,n} 	\\
			\rho_{2,1} & 1 & \cdots & \rho_{2,n} 	\\
			\vdots & \vdots &  & \vdots 				\\
			\rho_{n,1} & \rho_{n,2} & \cdots & 1
		\end{pmatrix}. 										\label{eq:korrelationsmatrix}
	\end{equation}
	Sowohl die Kovarianzmatrix $\Sigma$ als auch die Korrelationsmatrix $R$ sind per Definition quadratisch und auch symmetrisch. Diese Eigenschaften werden später noch nützlich sein.
	
	
	\subsection{Kondition einer Matrix}
	Ein zentrales Problem bei der Portfolio-Optimierung ist die numerische Stabilität. Diese wird maßgeblich durch die Konditionszahl $\kappa$ beeinflusst. Für eine invertierbare Matrix ist die Konditionszahl bezüglich der Spektralnorm definiert als das Verhältnis vom größten zum kleinsten Eigenwert:
	
	\begin{equation}
		\kappa(\Sigma) = \frac{\lambda_{\max}(\Sigma)}{\lambda_{\min}(\Sigma)} \quad \geq 1. \label{eq:kondition}
	\end{equation}
	
	Aus dieser Definition folgt direkt dass $\kappa$ immer größer oder gleich 1 sein muss. Denn für den Fall dass unsere Matrix nur einen Eigenwert hat erhalten wir für $\kappa$ den Wert 1. In allen anderen Fällen teilen wir den größten durch den kleinsten Eigenwert, was immer zu einem $\kappa$ $> 1$ führt.
	
	Die Konditionszahl lässt sich wie folgt interpretieren:
	
	Je kleiner die Konditionszahl, desto stabiler ist die Matrix. D.h. die Matrix reagiert weniger empfindlich gegenüber Änderungen in den Eingangsparametern.
	
	\clearpage
	
	\section{Das Markowitz-Modell}
	
	In den 50er Jahren des letzten Jahrhunderts entwickelte Harry Markowitz ein Modell für die klassische Portfoliotheorie. Das Ziel dieses Modells ist es, ein optimales Verhältnis zwischen dem Risiko und der Rendite herzustellen. Dieses Ziel soll durch Diversifikation erreicht werden.
	
	\subsection{Der rationale Investor}
	Markowitz geht für sein Modell von einem 'rationalem Investor' aus. Der rationale Investor verhält sich stets risikoavers und wird gleichzeitig jede Gelegenheit nutzen um Gewinne zu erzielen. Dieses Verhalten lässt sich durch ein $\mu$-$\sigma$-Diagramm veranschaulichen: \\
	
	%\vspace{1cm}
	\begin{figure}[htbp]
		\centering
		\includegraphics[width=0.5\textwidth]{Mu_Sigma_Kremer.pdf} 
		\caption{$\mu$-$\sigma$-Diagramm \parencite[11]{Kremer}}
		\label{fig:mu-sigma}
	\end{figure}
	
	Hat ein Investor die Wahl zwischen Investment $B$ und Investment $C$, so entscheidet er sich für $B$. Denn Investment $B$ hat die gleiche erwartete Rendite wie $C$, gleichzeitig aber ein geringeres Risiko.
	
	Muss er jedoch eine Auswahl zwischen Investment $A$ und Investment $C$ treffen, so entscheidet er sich für $A$. Investment $A$ liefert nämlich bei gleichem Risiko eine höhere erwartete Rendite.
	
	\clearpage
	
	\subsection{Erwartete Rendite und Varianz eines Portfolios}
	Hat man ein Portfolio bestehend aus $n$ Wertpapieren mit jeweils einem Gewicht $w$. Wobei die Summer der Gewichte 1 ergibt. So berechnet sich die erwartete Rendite des Portfolios als  \\
	
	\begin{equation}
		\mu_p = \sum_{i=1}^{n} w_i \mu_i = \langle w, \mu \rangle . \label{eq:portfoliorendite}
	\end{equation}
	
	Die Varianz eines Portfolios entspricht dem Risiko. Sie ergibt sich als \\
	
	\begin{equation}
		\sigma_p^2 = \sum_{i=1}^{n} \sum_{j=1}^{n} w_i w_j \sigma_{ij} = w^T \Sigma w . \label{eq:portfoliovarianz}
	\end{equation}
	
	
	\subsection{Diversifikation}
	Unter Diversifikation versteht man das Verteilen des Risikos auf mehrere Assets. Sie ist das fundamentale Prinzip hinter dem Modell von Markowitz. Die einfachste Form der Diversifikation ist die sogenannte \textbf{naive Diversifikation}. Hier werden einfach alle Assets gleich gewichtet. 
	Um uns die Diversifikationseffekte klar zu machen betrachen wir 2 Aktien $S^1$ und $S^2$ mit den Gewichten $w$ und $(1-w)$.
	
	Für die Portfoliorendite und die Portfoliovarianz gelten nach \textcite[22]{Kremer} die folgenden Zusammenhänge:
	
	\begin{equation}
		\mu_p = \mu_2 + w(\mu_1 - \mu_2) ,
	\end{equation}
	
	und
	
	\begin{equation}
		\sigma_p^2 = (w \sigma_1)^2 + ((1 - w) \sigma_2)^2 + 2w(1 - w)\sigma_1 \sigma_2 \rho .
	\end{equation}
	
	Wir wollen nun 2 Extremfälle betrachten um zu verstehen wie sich die Diversifikation auf ein Portfolio auswirkt:
	
	\textbf{$\rho = 1$}:
	
	Für $\rho = 1$ lässt sich die Portfoliovarianz schreiben als:
	
	\begin{equation*}
		\sigma_p^2 = (w \sigma_1)^2 + ((1 - w) \sigma_2)^2 + 2w(1 - w)\sigma_1 \sigma_2 .
	\end{equation*}
	
	Mit Hilfe der binomischen Formeln lässt sich das zu folgendem Ausdruck umformen:
	
	\begin{equation*}
		\sigma_p^2 = (w \sigma_1 + (1-w)\sigma_2)^2 ,
	\end{equation*}
	bzw.
	
	\begin{equation*}
		\sigma_p = \sigma_2 + w(\sigma_1 - \sigma_2) .
	\end{equation*}
	
	Wir sehen also, dass, wenn wir $w$ gleich 1 setzen, dass nur das Risiko von Wertpapier 1, also $\sigma_1$ übrig bleibt. Umgekehrt gilt das Gleiche für Wertpapier 2 bzw. $\sigma_2$. Für die erwartete Rendite gilt das Gleiche. D.h. für den Bereich zwischen $S^1$ und $S^2$ ergibt sich folgende Geradengleichung:
	
	\begin{equation*}
		\begin{pmatrix} \sigma_p \\ \mu_p \end{pmatrix} = \begin{pmatrix} \sigma_2 \\ \mu_2 \end{pmatrix} + w \begin{pmatrix} \sigma_1 - \sigma_2 \\ \mu_1 - \mu_2 \end{pmatrix} .
	\end{equation*}
	
	Wir sehen also, dass unsere Aktien $S^1$ und $S^2$ durch eine Gerade im $\mu$-$\sigma$-Diagramm verbunden sind. Wenn wir $w$ gleich 1 oder 0 setzen, landen wir genau auf $S^1$ bzw. $S^2$.
	
	\textbf{$\rho = -1$:}
	
	Setzen wir $\rho$ auf -1, dann erhalten wir:
	
	\begin{equation*}
		\sigma_p^2 = (w \sigma_1)^2 + ((1 - w) \sigma_2)^2 - 2w(1 - w)\sigma_1 \sigma_2 .
	\end{equation*}
	
	Auch hier kommen wir unter Verwendung der binomischen Formeln auf:
	
	\begin{equation*}
		\sigma_p^2 = (w \sigma_1 - (1-w)\sigma_2)^2
	\end{equation*}
	
	Da eine Standardabweichung nicht negativ sein kann, erhalten wir für $\sigma_p$
	
	\begin{equation*}
		\sigma_p = \left| w \sigma_1 - (1-w) \sigma_2 \right| .
	\end{equation*}
	
	Da in diesem Fall die Aktien negativ korreliert sind, kann man sich überlegen ob es ein Portfolio gibt, welches gar kein Risiko mehr aufweist. Hierzu setzen wir $\sigma_p$ gleich 0 und formen dann nach $w$ um. Dann erhalten wir
	
	\begin{equation*}
		w = \frac{\sigma_2}{\sigma_1 + \sigma_2} .
	\end{equation*}
	
	\clearpage
	
	 Für unser $\mu$-$\sigma$-Diagramm erhalten wir also folgende Fallunterscheidung:
	  
	\begin{equation*}
		\sigma_p = 
		\begin{cases} 
			\sigma_2 - w(\sigma_1 + \sigma_2) & \text{für } 0 \leq w \leq \frac{\sigma_2}{\sigma_1 + \sigma_2} \\
			w(\sigma_1 + \sigma_2) - \sigma_2 & \text{für } \frac{\sigma_2}{\sigma_1 + \sigma_2} < w \leq 1
		\end{cases} .
	\end{equation*}
	
	Wir sehen also, dass wir zwei Geraden ausgehend vom Nullpunkt auf der $\mu$-Achse erhalten. Für $-1 < \rho < 1$ liegen die Verbindungslinien also immer zwischen den beiden betrachteten Extremfällen $\rho = 1$ und $\rho = -1$. Im $\mu$-$\sigma$-Diagramm sieht das also folgendermaßen aus:
	
	 \begin{figure}[H]
	 \centering
	 \includegraphics[width=0.5\textwidth]{Rho.pdf} 
	 \caption{Diversifikation im $\mu$-$\sigma$-Diagramm \parencite[23]{Kremer}}
	 \label{fig:rho}
	 \end{figure}
	
	Wir sehen wo die verschiedenen Portfolios im $\mu$-$\sigma$-Diagramm liegen. In der Praxis treten der perfekt korrelierte und der perfekt unkorrelierte Fall allerdings nicht auf. Weshalb wir immer eine bogenförmige Verbindung zwischen zwei Aktien erhalten. Eine weitere Erkenntnis ist, dass es bei den bogenförmigen Kurven immer auch ein Portfolio gibt, wo das Risiko am geringsten ist. Nämlich den Punkt der im Diagramm am weitesten links liegt. Diesen Punkt erreicht man durch das Anpassen der Gewichte, also durch Diversifikation des Portfolios.
	
	\clearpage
	
	\subsection{Opportunitätsbereich}
	Da in der Realität immer $-1 < \rho < 1$ gilt. Lässt sich die Menge aller möglichen Portfolios im $\mu$-$\sigma$-Diagramm wie folgt darstellen.
	
	\begin{figure}[H]
	\centering
	\includegraphics[width=0.5\textwidth]{Oppertunitaetsbereich.pdf} 
	\caption{Opportunitätsbereich \parencite[29]{Kremer}}
	\label{fig:opportunität}
	\end{figure}
	
	
	Diese Menge von Portfolios bezeichnen wir als Opportunitätsbereich. Seine Form ähnelt einem "Regenschirm" \parencite[28]{Kremer}. Der Punkt $G$ bezeichnet das Portfolio mit dem global geringsten Risiko.
	

	\subsection{Effizienter Rand}
	Durch die Annahme des rationalen Investors können wir den Opportunitätsbereich verkleiner. Der rationale Investor investiert nur in die Portfolios die bei gleicher Rendite ein geringeres Risko aufweisen. Gleichzeitig fallen alle Portfolios weg, die bei gleichem Risiko eine geringere Rendite aufweisen. Die übrigen Portfolios liegen auf dem \textbf{Effizienten Rand} oder der \textbf{Effizienzlinie}.
	
	\begin{figure}[H]
	\centering
	\includegraphics[width=0.5\textwidth]{Effizienzlinie.pdf} 
	\caption{Effizienzlinie \parencite[30]{Kremer}}
	\label{fig:effizienter rand}
	\end{figure}
	
	Der effiziente Rand beschreibt also alle theoretisch optimalen Portfolios. In der Praxis stellt sich für einen Anleger die Frage nach dem spezifisch für ihn optimalen Portfolio. Diese Frage führt zu der Portfolio-Optimierung nach Markowitz.
	
	\subsection{Das Optimierungsverfahren}
	Nachdem der effiziente Rand die Menge aller theoretisch optimalen Portfolios definiert hat, besteht die praktische Aufgabe darin, das spezifisch für einen Investor optimale Portfolio auf dem effizienten Rand zu identifizieren. Dieser Prozess der Portfolio-Allokation wird als quadratisches Optimierungsproblem unter Nebenbedingungen formuliert.
	
	In dieser Arbeit liegt der Fokus auf der Risikominimierung, weshalb als Zielfunktion die vorher definierte Portfoliovarianz (vgl.Gl.\eqref{eq:portfoliovarianz}) dient
	
	\begin{equation}
		\min_w \sigma_p^2 = w^T \Sigma w . \label{eq:zielfunktion}
	\end{equation}
	
	\textbf{Bemerkung:}
	Alternativ könnte man als Zielfunktion einen Term wählen, der auch die Rendite enthält. Um bspw. Renditeziele vorzugeben und mit einem Risikoaversionparameter anzugeben, wie viel Risiko man bereit ist einzugehen. Diese Alternative Zielfunktion lautet
	
	\begin{equation}
		V(N) = \mu_p - \lambda \cdot \sigma_p^2 .
	\end{equation}
	
	\clearpage
	
	Das Optimierungsproblem \eqref{eq:zielfunktion} wäre trivialer Weise mit $w_i = 0$ gelöst. Da das Risiko am geringsten ist wenn wir in gar keine Aktie investieren, denn dann gilt $\sigma_p^2 = 0$. Weshalb man Nebenbedingungen formulieren muss.
	
	Die erste Nebenbedingungen ist die \textbf{Budgetrestriktion}. Hiermit stellt man sicher dass das gesamte Kapital investiert wird. Außerdem will man nicht über das Gesamtkapital hinaus gehen, also keine Kredite aufnehmen. Die erste Nebenbedingung lautet also
	
	\begin{equation}
		\sum_{i=1}^{n} w_i = 1 . \label{eq:budgetrestriktion}
	\end{equation}
	
	Neben der Budgetrestriktion wird in dieser Arbeit auch eine Nebenbedingung die Leerverkäufe ausschließt benötigt. Hierfür nutzt man eine \textbf{Nicht-Negativitäts-Bedingung}. Es soll also stets gelten
	
	\begin{equation}
		w_i \geq 0 \quad \text{für alle } i = 1, ..., n \label{eq:nicht-negativ}
	\end{equation}
	
	Da die Kovarianzmatrix $\Sigma$ per Definition positiv definit ist, handelt es sich um ein konvexes Optimierungsproblem mit linearen Nebenbedingungen. Für solche Probleme existiert eine eindeutige Lösung. Das aus der Lösung resultierende Portfolio nennen wir \textbf{Minimum-Varianz-Portfolio} (MVP).
	
	Um die mathematische Struktur zu verdeutlichen, betrachten wir zunächst das Optimierungsproblem nur unter der Budgetrestriktion. Dieses Problem lässt sich analytisch mit der Lagrange-Funktion lösen.
	
	Die Lagrange-Funktion hat die Struktur
	
	\begin{equation}
		L(w,\lambda) = \operatorname{Zielfunktion} - \lambda \cdot (\operatorname{Nebenbedingung} = 0) . \label{eq:lagrange}
	\end{equation}
	
	Die Zielfunktion wurde in \eqref{eq:zielfunktion} bereits definiert. Die Budgetrestriktion nach null umgestellt lautet $\sum_{i=1}^{n} w_i - 1$. Also erhält man für die Lagrange-Funktion den Ausdruck
	
	\begin{equation*}
		L(w,\lambda) = w^T \Sigma w - \lambda (\mathbf{1}^T w-1) .
	\end{equation*}
	
	Die Nebenbedingung wurde hier von der Summenschreibweise in die Vektorschreibweise übertragen. Um das Minimum zu bestimmen, müssen die partiellen Ableitungen nach $w$ und nach $\lambda$ gebildet und gleich Null gesetzt werden:
	
	\begin{equation}
		\frac{\partial L}{\partial w} = 2 \Sigma w - \lambda \mathbf{1} \stackrel{!}{=} 0 \label{eq:abl_w}
	\end{equation}
	und
	\begin{equation}
		\frac{\partial L}{\partial \lambda} = \mathbf{1}^T w - 1 \stackrel{!}{=} 0. \label{eq:abl_lambda}
	\end{equation}
	
	Aus Gleichung \eqref{eq:abl_w} lässt sich nun der Zusammenhang zwischen den optimalen Gewichten und der Kovarianzmatrix isolieren. Durch Umstellen ergibt sich:
	
	\begin{equation}
		w = \frac{1}{2} \lambda \Sigma^{-1} \mathbf{1}. \label{eq:w_struktur}
	\end{equation}
	
	Um den Lagrange-Multiplikator $\lambda$ zu eliminieren setzen wir \eqref{eq:w_struktur} in die Budgetrestriktion \eqref{eq:budgetrestriktion} ein. Dadurch erhalten wir den Ausdruck
	
	\begin{equation*}
		\mathbf{1}^T \left(\frac{1}{2} \lambda \Sigma^{-1} \mathbf{1} \right) = 1.
	\end{equation*}
	
	$\frac{1}{2} \lambda$ ist eine Skalar, weshalb man es in der Gleichung nach vorne ziehen kann. Stellen wir die gesamten Gleichung jetzt nach genau diesem Skalar um, so erhalten wir
	
	\begin{equation*}
		\frac{1}{2} \lambda = \frac{1}{\mathbf{1}^T \Sigma^{-1} \mathbf{1}}
	\end{equation*}
	
	Setzt man diesen Term nun in die Gleichung \eqref{eq:w_struktur} ein, erhält man für die Gewichte des Minimum-Varianz-Portfolios die Gleichung
	
	\begin{equation}
		w_{MVP} = \frac{\Sigma^{-1}\mathbf{1}}{\mathbf{1}^T \Sigma^{-1} \mathbf{1}} . \label{eq:w_mvp}
	\end{equation}
	
	% Evtl. in diesem Satz nochmal das genau Kapitel erwähnen, wo die Bezugnahme statt findet.
	
	Man sieht also, dass die Gewichte des Minimum-Varianz-Portfolios direkt von der Inversen der Kovarianzmatrix abhängen. Auf diese Eigenschaft wird in einem späteren Kapitel noch einmal Bezug genommen werden.
	
	\textbf{KKT-Bedingungen für das Portfolioproblem:}	
	
	Das Optimierungsverfahren wie es bisher betrachtet wurde, berücksichtigt nur die Budgetrestriktion. Allerdings möchten wir in dieser Arbeit Leerverkäufe ausschließen. Nimmt man also diese Nebenbedingung hinzu, kann man nicht mehr so einfach eine analytische Lösungen für das Minimum-Varianz-Portfolio bestimmen. Denn es existiert dann zusätzlich noch eine Ungleichheitsnebenbedingung \eqref{eq:nicht-negativ}. Um dieses Optimierungsproblem zu lösen wird die Langrange-Gleichung um die der Kuhn-Tucker-Bedinungen erweitert.
	
	Allgemein erhält man für die Langrange-Gleichung mit Kuhn-Tucker den Ausdruck \parencite[260]{Nickel}
	
	\begin{equation}
		L(w, \lambda, \delta) = f(w) + \sum_{i=1}^{l} \lambda_i \cdot h_i(w) + \sum_{j=1}^{m} \delta_j \cdot g_j(w) . \label{eq:langrange_kkt}
	\end{equation}
	
	$h$ beschreibt in dieser Gleichung die nach null umgestellten Gleichheitsnebenbedingungen und $g$ beschreibt die nach null umgestellten Ungleichheitsnebenbediungen.
	
	
	Um nun den Vektor der optimalen Portfoliogewichte zu bestimmen, müssen die folgenden Kuhn-Tucker-Bedingungen erfüllt sein \parencite[260]{Nickel}
	
	\begin{equation}
		\nabla L = \nabla f(w^*) + \sum_{i=1}^{l} \lambda_i \nabla h_i(w^*) + \sum_{j=1}^{m} \delta_j \nabla g_j (w^*) = 0 .
	\end{equation}
	
	$w^*$ bezeichnet den Vektor mit den optimalen Portfoliogewichten. Die weiteren Kuhn-Tucker-Bedingungen lauten
	
	\begin{enumerate}
		\item $h_i(w^*) = 0$
		\item $g_j(w^*) \leq 0$
		\item $\delta_j \geq 0$
		\item $\delta_j \cdot g_j(w^*) = 0 \quad \forall j=1, ..., m$ .
	\end{enumerate}
	
	Die erweiterte Langrange-Gleichung für das in dieser Arbeit betrachtete Minimum-Varianz-Portfolio lautet also
	
	\begin{equation}
		L = w^T \Sigma w + \lambda (\mathbf{1}^Tw - 1) + \sum_{j=1}^{n} \delta_j  (-w_j) . \label{eq:kkt_mvp}
	\end{equation}
	
	\textbf{Numerische Lösung des Optimierungsproblems: }
	
	Ein solches Optimierungsproblem kann man numerisch mithilfe des Critical-Line-Algorithmus (CLA) lösen. Die Idee des CLA besteht darin, den effizienten Rand nicht durch unzählige Einzelberechnungen zu schätzen, sondern ihn systematisch durch die Identifikation von sogenannten Eckportfolios (Corner Portfolios) exakt zu konstruieren. Diese Eckportfolios sind genau die Punkte, an denen entweder eine Aktie die Nullschranke erreicht und aus dem Portfolio ausscheidet oder eine Aktie neu in das Portfolio aufgenommen wird. Durch dieses sukzessive Abwandern des effizienten Randes liefert der CLA eine exakte Lösung für das Minimum-Varianz-Portfolio und berücksichtigt dabei alle Ungleichheitsrestriktionen.
	
	\subsection{Kritik am Markowitz-Modell}
	
	Obwohl der Critical-Line-Algorithmus das Minimum-Varianz-Portfolio exakt berechnet, gibt es berechtigte Kritik am Modell von Markowitz. Die Probleme liegen dabei nicht in der Rechenmethode selbst, sondern in der Beschaffenheit der Eingangsdaten. 
	
	Es wurde bereits in Gleichung \eqref{eq:w_mvp} gezeigt, dass die Portfoliogewichte direkt mit der Inversen der Kovarianzmatrix zusammenhängen. Da diese Matrix allerdings, wie wir später noch sehen werden, numerisch sehr instabil ist, führt eine kleine Fehlschätzung der Varianzen bzw. Kovarianzen zu großen Veränderungen im Gewichtsvektor. Da alle Einträge der Kovarianzmatrix aus historischen Daten geschätzt werden, kann dies zu Problemen führen.
	
	Ein weiteres Problem ist, dass Markowitz dazu neigt die Gewichte zu "konzentrieren". Das bedeutet, dass wenige Aktien stark gewichtet werden, während viele andere Aktien mit einem Gewicht von null versehen werden. Das Modell neigt dazu hoch volatile Aktien mit null zu gewichten, was natürlich im ersten Moment logisch erscheint, denn man will ja eigentlich ein Portfolio mit möglichst geringer Volatilität (MVP). Doch dieses Verhalten steht im Konflikt mit dem Diversifikationsgedanken. Denn das Ziel eines Anlegers ist es das Risiko auf mehrere Aktien zu verteilen. 
	
	Der letzte Kritikpunkt am Markowitz-Modell, der hier betrachtet wird, ist, dass dieses Modell die Kovarianzen der Vergangenheit betrachtet. Oft sind die Kovarianzen zwischen zwei Aktien aber keine konstanten Werte, sondern sie verändern sich mit der Zeit. D.h. die Annahme dass die Kovarianzen zwischen den Aktien immer den Kovarianzen des vorherigen Zeitraums entsprechen ist in der Praxis häufig falsch. Hierbei handelt es sich aber um eine notwendige Modellannahme, denn man kann die Kovarianzen der Zukunft nicht vorhersagen.
	
	Zusammenfassend lässt sich festhalten, dass die klassische Portfolio-Optimierung nach Markowitz trotz ihrer mathematischen Eleganz signifikante Schwächen in der Praxis aufweist. Weshalb wir uns im nächsten Kapitel mit einem alternativen Modell beschäftigen, dem \textbf{Hierachical-Risk-Parity-Modell}.
	
	\clearpage
	
	\section{Das Hierarchical-Risk-Parity-Modell}
	
	Um die im letzten Kapitel angesprochenen Schwächen der klassischen Portfolio-Optimierung zu umgehen, hat Marcos López de Prado im Jahr 2016 das Hierarchical-Risk-Parity-Modell entwickelt. Dieses Modell verzichtet komplett auf die Matrixinversionen. Stattdessen nutzt HRP Methoden des maschinellen Lernens, um die hierarchische Korrelationsstruktur der Assets abzubilden und für die Gewichtung zu nutzen.
	
	Der HRP-Agorithmus lässt sich in drei Phasen unterteilen:
	
	\begin{enumerate}
		\item Das \textbf{hierarchische Clustering}, um die Ähnlichkeitsstruktur der Assets zu bestimmen
		\item Die \textbf{Quasi-Diagonalisierung}, um die Kovarianzmatrix basierend auf dieser Struktur neu zu ordnen
		\item Die \textbf{Recursive Bisection}, um das Kapital rekursiv über die Hierarchieebenen zu verteilen.
	\end{enumerate}
	
	\subsection{Hierarchisches Clustering} % Linkage vielleicht noch erwähnen, bzw. genauer darauf eingehen wie dieses Clustering funktioniert (evtl.)
	
	Der erste Schritt des Modells besteht darin die Assets in einer hierarchischen Baumstruktur anzuordnen. Die Idee hierbei ist, dass die Assets keine homogene Masse sind, sondern es eine natürliche Gruppierung von Assets gibt die sich ähnlich verhalten (bspw. alle Aktien der Automobilbranche).
	
	Um diese Struktur mathematisch zu erfassen wird zunächst die Korrelationsmatrix $R$ der Renditen betrachtet \eqref{eq:korrelationsmatrix}. Aus dieser Korrelationsmatrix leitet man die Distanzmatrix ab. Hierfür nutzt man die Metrik \parencite[224]{Lopez}
	
	\begin{equation}
		d_{i,j} = \sqrt{\frac{1}{2} (1- \rho_{i,j})} . \label{eq:dist_metrik}
	\end{equation}
	
	Diese Distanzmetrik sorgt dafür dass perfekt positiv korrrelierte Assets ($\rho = 1$) eine Distanz von null aufweisen und perfekt negativ korrelierte Assets ($\rho = -1$) eine Distanz von 1. Auf Basis dieser Distanzen wird eine agglomeratives Clusteringverfahren angewendet, dabei werden die Assets sukzessive zu Paaren und zu Clustern zusammengefasst. Das Ergebnis ist ein Dendrogramm, welches die Nähe der verschiedenen Assets zueinander veranschaulicht.
	
	% Für die Erstellung des Dendrogramms muss zudem ein Linkage-Verfahren gewählt werden, welches definiert, wie die Distanz zwischen neu entstehenden Clustern berechnet wird. In dieser Arbeit wird das Ward-Verfahren angewendet. Dieses minimiert die Zunahme der Varianz innerhalb der Cluster bei jedem Fusionsschritt und führt tendenziell zu einer sehr ausgeglichenen und hierarchisch schlüssigen Baumstruktur, was für die spätere Risikoverteilung vorteilhaft ist.
	
	\begin{figure}[H]
		\centering
		\includegraphics[width=0.7\textwidth]{dendrogramm.pdf} 
		\caption{Dendrogramm für 5 Aktien}
		\label{fig:dendrogramm_5}
	\end{figure}
	
	Wir sehen hier ein beispielhaftes Dendrogramm, bestehend aus den Aktien von BMW, Mercedes-Benz, der Allianzversicherung, BASF und Bayer. Man erkennt die 2 Hauptcluster, Cluster 1 besteht aus den beiden Aktien von BMW und Mercedes-Benz und Cluster 2 beinhaltet die Aktien von BASF und Bayer. Die Allianz bildet ein eigenes Cluster. Man kann also direkt sehen, dass die Aktien nach Branchen geclustert wurden. Auf der Y-Achse sind die Distanzen verzeichnet. Die Automobilaktien liegen ca. bei einer Distanz von 0,4 während die Aktien der Pharmaindustrie etwas darunter liegen (ca. 0,37). Die Allianz erhält ihr eigenes Cluster, da ihr Abstand zu den anderen Aktien höher ist, sie unterscheidet sich also zu stark vom Rest. Die Distanz zu den Aktien von BASF und Bayer liegt ca. bei 0,45. 
	
	Nach dem durch das Clustering die hierarchische Struktur der Assets identifiziert wurde, liefert die Anordnung der Endknoten die Sortierung für den zweiten Schritt des HRP-Algorithmus. In Abbildung \ref{fig:dendrogramm_5} ergibt sich z.B. von links nach rechts die Reihenfolge: BMW, Mercedes, Allianz, BASF, Bayer. 
	
	\subsection{Quasi-Diagonalisierung}
	
	Die aus dem Clustering gewonnene Sequenz nutzt man jetzt für die Quasi-Diagonalisierung. Bei diesem Schritt geht es darum die Zeilen und Spalten der ursprünglichen Korrelationsmatrix so anzuordnen, dass stark korrelierende Aktien (bzw. Aktien mit einer geringen Distanz) in der Matrix nebeneinander angeordnet sind. Durch diese neue Ordnung konzentrieren sich die hohen Korrelationen um die Hauptdiagonale der Matrix. Weshalb man dieses Vorgehen als \textbf{Quasi-Diagonalisierung} bezeichnet. Im Gegensatz zur Markowitz-Optimierung, ist in diesem Schritt keine Matrixinversion nötig. Das macht diesen Schritt numerisch stabiler.
	% Im Gegensatz zur klassischen Markowitz-Optimierung ist für diesen Schritt keine Inversion der Matrix notwendig, was das Modell besonders robust gegenüber Multikollinearität macht
	
	\subsection{Recursive Bisection}
	
	Der finale Schritt des HRP-Algorithmus ist die \textbf{Recursive Bisection}. Dabei wird das Kapital auf Basis der quasi-diagonalen Struktur schrittweise von oben nach unten (Top Down) verteilt. Im Gegensatz zu Markowitz, wo die Gewichte simultan bestimmt werden, arbeitet der HRP-Algorithmus rekursiv entlang der Baum-Hierarchie.
	
	\textbf{Bestimmung der Gewichte: }
	
	Das Portfolio wird zunächst in zwei Cluster unterteilt. Für die zwei Teilcluster $C_1$ und $C_2$ wird die Varianz bestimmt. Dabei wird innerhalb der Cluster eine \textbf{Inverse-Variance-Gewichtung} (IVP) unterstellt:
	
	\begin{equation}
		Var(C_i) = g_i^T \Sigma_i g_i . \label{eq:ivp}
	\end{equation}
	
	Wobei $g_i$ die Gewichtung innerhalb des Clusters bezeichnet. Die einzelnen Komponenten berechnen sich als
	
	\begin{equation}
		g_k = \frac{\frac{1}{\sigma_k^2}}{\sum_{j \in C_i} \frac{1}{\sigma_j^2}} . \label{eq:ivp_weights}
	\end{equation}
	
	Um also die Komponenten für $g_i$ zu bestimmen, nimmt man den Kehrwert der Varianz des k-ten Assets und teilen diesen Wert durch die Summe der Kehrwerte aller Assets, die sich in diesem Cluster befinden.
	
	Die Gewichte zwischen den Clustern berechnet man als
	
	\begin{equation}
		\alpha = 1 - \frac{Var(C_1)}{Var(C_1) + Var(C_2)} . 
	\end{equation}
	
	Das Cluster $C_1$ erhält den Anteil $\alpha$ während das Cluster $C_2$ den Anteil $(1 - \alpha)$ erhält. Diesen Prozess wiederholt man so lange bis jedes Asset ein finales Gewicht zugeordnet bekommen hat.
	
	Dieses rekursive Verfahren kann man sich nun Anhand des in Abbildung \ref{fig:dendrogramm_5} vorgestellten Beispiels mit den fünf DAX-Aktien verdeutlichen.
	
	\clearpage
	
	% Start Gemini Text
	
	Basierend auf der im Dendrogramm (Abbildung \ref{fig:dendrogramm_5}) sichtbaren Hierarchie und der durch die Quasi-Diagonalisierung ermittelten Reihenfolge [BMW, Mercedes, Allianz, BASF, Bayer] vollzieht der Algorithmus die Allokation in folgenden Schritten:
	
	\begin{enumerate}
		\item \textbf{Erste Bisektion (Top-Level):} Das Gesamtportfolio wird zunächst in zwei Hauptgruppen unterteilt. In diesem Fall bildet der Automobil-Sektor (BMW und Mercedes) das erste Cluster $C_1$, während die restlichen drei Assets (Allianz, BASF und Bayer) das Cluster $C_2$ bilden. Gemäß der Varianz dieser beiden Blöcke wird das Gesamtkapital erstmals aufgeteilt.
		\item \textbf{Zweite Bisektion:} Innerhalb des zweiten Clusters findet eine weitere Aufteilung statt. Hierbei wird die Allianz als einzelnes Asset von dem verbleibenden Paar aus BASF und Bayer getrennt. Das Gewicht von $C_2$ wird also zwischen der Allianzversicherung und dem Chemie-Cluster aufgeteilt.
		\item \textbf{Finale Ebene:} In der untersten Hierarchiestufe wird das Kapital innerhalb der eng verwandten Paare verteilt. Das für den Automobil-Sektor reservierte Kapital wird zwischen BMW und Mercedes aufgeteilt; analog wird das Kapital des Chemie-Clusters zwischen BASF und Bayer partitioniert.
	\end{enumerate}
	
	Diese Vorgehensweise verdeutlicht den entscheidenden Vorteil der {Recursive Bisection}: Das Risiko wird zunächst zwischen ökonomisch unterschiedlichen Sektoren (Automobil vs. Finanzen/Chemie) gestreut. Erst innerhalb dieser stabilen Blöcke findet die Feinverteilung auf die Einzelwerte statt. Dies verhindert, dass hohe Korrelationen zwischen ähnlichen Aktien – wie etwa BMW und Mercedes – zu einer instabilen Übergewichtung eines einzelnen Titels führen, wie es bei der klassischen Rendite-Varianz-Optimierung oft der Fall ist.
	
	% Ende Gemini Text
	
	\subsection{Fazit zu HRP}
	
	Zusammenfassend kann man festhalten, dass das Hierarchical-Risk-Parity-Modell von Marcos López de Prado eine Brücke zwischen der klassischen Portfoliotheorie und dem maschinellen Lernen schlägt. Durch den Verzicht auf die Inversion von Matrizen und die Konzentration auf die hierarchische Struktur wird eine höhere Robustheit gegen Schätzfehler in den Inputparametern angestrebt. Ob sich diese theoretischen Vorteilen in der Praxis bewähren, wird im folgenden Kapitel mit einem Portfolio aus 30 Aktien empirisch untersucht. Dabei wird insbesondere die Konditionszahl als Maß für die numerische Stabilität im Vergleich zur klassischen Optimierung herangezogen.
	
	\clearpage
	
	\section{Empirischer Vergleich}
	
	Nachdem die theoretischen Grundlagen sowohl für das Markowitz-Modell als auch für das Hierarchical-Risk-Paritiy Modell geklärt wurden, folgt im vorliegenden Kapitel der empirische Vergleich der beiden Modelle. Hierfür werden beide Modelle auf ein Portfolio bestehend aus 30 Assets angewendet. Als Benchmark dient in diesem Vergleich die naive Diversifikation. Also ein Portfolio in dem alle Assets gleich, also mit $\frac{1}{N}$, gewichtet werden. Der Zeitraum den wir betrachten wollen erstreckt sich über 20 Jahre (2006 - 2026).
	
	\subsection{Zusammenstellung des Portfolios}
	
	Die Auswahl der 30 Assets für dem empirischen Vergleich folgt einer gezielten Strategie, um die Robustheit der Optimierungsalgorithmen unter verschiedenen Marktbedingungen und Korrelationsstrukturen zu testen. Das Portfolio wurde nach folgenden Kriterien zusammengestellt:
	
	\begin{enumerate}
		\item \textbf{Geografische Diversifikation:} Um das Klumpenrisiko einzelner Volkswirtschaften zu reduzieren, deckt das Portfolio mehrere Wirtschaftsregionen ab.
		
		- \textbf{Deutschland:} Die 10 DAX-Aktien decken eine Volkswirtschaft in der Eurozone ab. Vor allem die Industrie und den Finanzsektor.
		
		- \textbf{USA:} Die zehn US-Titel decken primär den Technologie- und Konsumgütersektor ab, welcher historisch andere Wachstums- und Volatilitätsprofile aufweist als europäische Werte.
		
		- \textbf{Japan:} Die neun japanischen Titel bringen eine zusätzliche Diversifikation durch die geringere Korrelation des japanischen Marktes zu westlichen Märkten ein.
		
		\item \textbf{Abdeckung von Risikofaktoren:} Die Aufnahme des MSCI Emerging Markets dient als Proxy für die Märkte in Entwicklungsländern. Da sich diese Märkte oft anders verhalten als die Ökonomien von Industrieländern, erhöht dieser Titel nochmal die Heterogenität des Portfolios. 
		
		\item \textbf{Datenverfügbarkeit:} Es wurden ausschließlich Titel gewählt, die im betrachteten Zeitraum (2006 - 2026) eine kontinuierliche Datenbasis liefern. Das stellt sicher, dass die betrachteten Kovarianzmatrizen nicht durch fehlende Daten verzerrt werden.
		
		\item \textbf{Sektorale Streuung:} Innerhalb der verschiedenen Regionen wurde darauf geachtet, dass die Assets aus verschiedenen Sektoren (Industrie, Finanzen, Pharma, Technologie, Konsum) stammen. Da HRP die Aktien nach den Sektoren clustert, ist es wichtig verschiedenen Sektoren abzubilden, um den potenziellen Vorteil dieser Strategie überhaupt sichtbar zu machen.
		
	\end{enumerate}
	
	\clearpage
	
	Folgende Assets wurden nach den oben genannten Kriterien ausgewählt:
	
	\begin{table}[htbp]
		\centering
		\caption{Zusammensetzung des Portfolios nach Regionen und Unternehmen}
		\label{tab:portfolio_names}
		\small % Verkleinert die Schrift etwas, damit alles sauber passt
		\begin{tabularx}{\textwidth}{|l|X|}
			\hline
			\textbf{Region} & \textbf{Unternehmen / Assets} \\ \hline
			\textbf{Deutschland} & SAP, Siemens, Allianz, Münchener Rück, Deutsche Telekom, DHL Group, Rheinmetall, E.ON, Infineon, Beiersdorf \\ \hline
			\textbf{USA} & NVIDIA, Microsoft, Alphabet, Amazon, UnitedHealth, Visa, Procter \& Gamble, Costco, NextEra Energy, Lockheed Martin \\ \hline
			\textbf{Japan} & Toyota, Sony, Tokyo Electron, Keyence, Mitsubishi UFJ, Nintendo, Shin-Etsu Chemical, Takeda Pharmaceutical, Fast Retailing, Itochu \\ \hline
			\textbf{Global} & MSCI Emerging Markets ETF (EEM) \\ \hline
		\end{tabularx}
		\vspace{5pt}
		\par\textit{Quelle: Eigene Darstellung, Kursdaten von Yahoo Finance.}
	\end{table}
	
	\subsection{Methodik und Parameter}
	
	Um eine belastbare Vergleichbarkeit der Portfoliomodelle im empirischen Kontext zu gewährleisten, ist die Definition einheitlicher Rahmenbedingungen und Modellparameter erforderlich. Der Vergleich erfolgt als rollierender Backtest über den Zeitraum von 2006 bis 2026. Die Wahl des Zeitraums wurde deshalb so getroffen, da sowohl die Finanzkrise 2008 als auch die Covid-19 Pandemie innerhalb dieses Zeitraums liegen.
	
	Für die Durchführung der Simulation wurden folgende Kernparameter festgelegt:
	
	\begin{itemize}
		\item \textbf{Schätzzeitraum (Lookback):} Zur Berechnung der Eingabeparameter wird ein rollierendes Fenster von 12 Monaten verwendet. Dieser Zeitraum stellt sicher dass man genügend Datengrundlage für das Rebalancing hat, jedoch auch auf aktuelle Trends im Markt reagieren kann. Als Datengrundlage dienen die Schlusskurse der einzelnen Handelstage.
		\item \textbf{Rebalancing-Frequenz:} Die Portfolios werden jährlich, halbjährlich sowie quartalsweise neu gewichtet. Die Rebalancing Frequenzen orientieren sich an realistischen Zeiträumen aus der Finanzwelt. Es werden drei verschiedenen Frequenzen betrachtet, um zu prüfen ob sich die Modelle abhängig vom Rebalancing unterschiedlich verhalten, oder ob die Frequenz keine merklichen Auswirkungen auf die jeweilige Performance hat.
		\item \textbf{Benchmark:} Als Referenz dient die naive Diversifikation (1/N). Das Benchmarkportfolio dient als Vergleichswert. Ziel dabei ist die Überprüfung wie groß der Optimierungseffekt auf die Portfolios mit den verschiedenen Ansätzen überhaupt ist. 
	\end{itemize}
	
	\clearpage
	
	Um die zeitliche Additivität der Renditen zu gewährleisten und die numerische Stabilität der Schätzungen zu verbessern, werden Log-Renditen verwendet. Sie berechnen sich als
	
	\begin{equation}
		r_{t} = \ln \left( \frac{S_t}{S_{t-1}} \right) . \label{eq:log_returns}
	\end{equation}
	
	
	
	Wie bereits in Kapitel 3.7 dargelegt, stellt die numerische Instabilität bei der Inversion der Kovarianzmatrix eine zentrale Schwachstelle des Markowitz-Modells dar. Im Rahmen dieser empirischen Untersuchung wird daher die Konditionszahl \eqref{eq:kondition} der Kovarianzmatrizen explizit überwacht. Dies erlaubt es, die theoretisch postulierten Vorteile des HRP-Algorithmus – insbesondere den Verzicht auf die Matrixinversion – anhand der vorliegenden Marktdaten zu validieren.
	
	Für die praktische Umsetzung des klassischen Markowitz-Modells wurde im Optimierungsschritt das sogenannte Shrinkage-Verfahren eingesetzt. Hierbei werden kleine Werte auf der Diagonalen der Kovarianzmatrix addiert, da die Matrix sonst zu nah an der Singularität wäre. Das würde dazu führen dass das Optimierungsverfahren bei gewissen Rebalancing-Schritten nicht konvergiert. Durch diese Regularisierung wird die numerische Stabilität des Algorithmus verbessert, um eine robuste Schätzung der Portfoliogewichte sicherzustellen.
	Die regularisierten Kovarianzmatrizen berechnen sich also als
	
	\begin{equation}
		\Sigma_{i, \epsilon} = \Sigma_i + I \cdot \epsilon . \label{eq:shrinkage}
	\end{equation}
	
	Hierbei stellt $I$ die Einheitsmatrix dar und $\epsilon$ ist ein sehr kleiner Wert größer Null.
	
	Als weiteres Bewertungskriterium für die Performance der verschiedenen Modelle, wird das Sharpe-Ratio betrachtet. Es setzt die Überrendite des Portfolios $\mu_p$ ins Verhältnis zu seiner Volatilität $\sigma_p$. Unter der Annahme eines riskofreien Zinssatzes $r_f$ berechnet sich das Sharpe-Ratio wie folgt:
	
	\begin{equation}
			SR = \frac{\mu_p - r_f}{\sigma_p} \label{eq:sharpe_ratio}
	\end{equation}
	
	Das Sharpe-Ratio misst das Verhältnis von Rendite zur Volatilität. Ein hohes Sharpe-Ratio deutet also auf eine effizientere Risikoallokation hin, während ein niedriges Sharpe-Ratio das Gegenteil andeutet. Es dient als zentrales Maß für die Diversifikation eines Portfolios.
	
	Da im betrachteten Zeitraum über einen Großteil der Zeit eine sehr starke Niedrigzinspolitik herrschte, wird in dieser Untersuchung der risikolose Zins $r_f$ mit null angenommen. Außerdem liegt der Fokus auf dem Vergleich der Modelle. Da der risikolose Zins bei allen Portfolios gleichermaßen abgezogen werden würde, beeinträchtigt diese Vereinfachung den Vergleich nicht.
	
	\clearpage
	
	\subsection{Ergebnisse und Vergleich}
	
	Zunächst wird die numerische Stabilität betrachtet, denn das ist einer der Hauptkritikpunkte von Marcos López de Prado am Markowitz-Modell \parencite[222]{Lopez}.
	
	Da die Kovarianzmatrizen alle symmetrisch sind (siehe \eqref{eq:kovarianzmatrix}), wurde die Konditionszahl in dieser Analyse mit \eqref{eq:kondition} bestimmt. Plottet man nun die Konditionszahl gegen die Zeit erhält man folgenden Graphen.
	
	% Beim Konditions-Plot die Markierung bei 10³ entfernen
	
	\begin{figure}[H]
	\centering
	\includegraphics[width=1.0\textwidth]{kondition_hj.pdf} 
	\caption{Konditionszahl (halbjährlich) über die Zeit}
	\label{fig:kondition}
	\end{figure}
	
	% evtl die anderen Plots zur Konditionszahl im Anhang zeigen
	% Quelle für kritsche Konditionszaheln (200, 300)
	Exemplarisch wird hier die rohe Konditionszahl bei halbjährlicher Rebalancing Frequenz betrachtet, d.h. es wird die Konditionszahl berechnet, bevor das Shrinkage-Verfahren angewendet wird. Man kann sehen, dass die Konditonszahl über die Zeit stark schwankt. Jedoch erreicht sie nie den kritischen Wert von $10^3$. Allerdings gilt auch eine Konditionszahl von einem Wert über 200 bereits als numerisch instabil. In diesem Bereich bewegt sich die Konditionszahl der Kovarianzmatrizen in regelmäßigen Abständen. Dieser Umstand erklärt auch, wieso die Anwendung von Shrinkage überhaupt notwendig ist, denn gewisse Matrizen sind einfach zu schlecht Konditioniert für eine stabile Schätzung der Gewichte.
	
	\begin{table}[H]
	\centering
	\caption{Analyse der numerischen Stabilität (Konditionszahl $\kappa$)}
	\label{tab:stabilitaet_vergleich}
	\begin{tabular}{@{}lccc@{}}
		\toprule
		\textbf{Rebalancing-Frequenz} & \textbf{Max. $\kappa$} & \textbf{Ø $\kappa$} & \textbf{Stabilitätsrisiko} \\ \midrule
		Quartalsweise & 229,33 & 103,42 & Hoch \\
		Halbjährlich & 225,59 & 103,49 & Hoch \\
		Jährlich & 221,69 & 103,47 & Hoch \\ \bottomrule
	\end{tabular}
	\end{table}
	
	Wenn man die Konditionszahl für verschiedene Rebalancing Frequenzen betrachtet, stellt man fest dass die durchschnittliche Konditionszahl fast konstant ist. Sie ändert sich lediglich an der zweiten Nachkommastelle. Die maximale Konditionszahl nimmt mit niedriger Rebalancing-Frequenz ab, jedoch auch nur geringfügig. Die Rebalancing-Frequenz scheint nur einen sehr geringen Einfluss auf die Konditionierung der Kovarianzmatrizen zu haben.
	
	Diese Untersuchung bestätigt also die Behauptung von Marcos López de Prado, dass die Matrizen für die Markowitz-Optimierung schlecht konditioniert sein können (vgl. \textcite[222]{Lopez}), und dadurch kleine Fehler in den geschätzten Parametern zu großen Veränderungen in den berechneten Portfoliogewichten führen können. Es wird auch klar, warum die es sinnvoll sein kann einen Algorithmus zu verwenden, der auf Matrixinversionen verzichtet, und somit von numerischen Instabilitäten deutlich weniger betroffen ist.
	
	
	\textbf{Performance:}
	
	Um die Performance der Modelle zu bewerten, werden Kennzahlen wie die Rendite, Volatilität, das Sharpe-Ratio und die Verteilung der Gewichte betrachtet.
	
	Als erstes betrachtet man die allgemeine Performance. Hierzu schaut man sich die Rendite der verschiedenen Anlagestrategien über die Zeit an. Exemplarisch wird hier die Rendite über die Zeit bei halbjährlicher Rebalancing-Frequenz betrachtet.
	
	\begin{figure}[H]
	\centering
	\includegraphics[width=1.0\textwidth]{performance_hj.pdf} 
	\caption{Performance der Modelle bei halbjährlicher Rebalancing-Frequenz}
	\label{fig:performance}
	\end{figure}
	
	Im Performance-Plot der Modelle stellt man fest, dass die naive Diversifikation die höchste Rendite erzielt. Auch wenn Markowitz und HRP im Zeitraum von 2012 - 2016 höhere Renditen erzielen, werden sie am Ende deutlich von dem $\frac{1}{N}$-Portfolio überholt.
	
	Die Graphen der verschiedenen Portfolios verlaufen ziemlich parallel. Weshalb man gut sieht, dass die Ausschläge sowohl bei HRP als auch bei Markowitz weniger stark sind als bei der naiven Diversifikation. Daraus kann man schlussfolgern dass die Allokations-Modelle die Volatilität, also das Risiko, tatsächlich reduziert haben. Der Preis für dieses reduzierte Risiko ist demnach die verringerte Rendite. Außerdem sieht man ca. im Jahr 2021 also zur Covid-Krise den vermutlich größten Wertverlust (Max. Drawdown) der Portfolios. Hier scheinen sowohl das HRP-Modell, als auch Markowitz einen geringeren Wertverlust zu erleiden als die naive Diversifikation.
	
	Um die Gewichtsverteilung der Modelle zu bewerten werden hier Area-Plots betrachtet. Sie zeigen wie sich die Gewichte über die Zeit verändert haben. Die naive Diversifikation wird hierbei nicht betrachtet, denn dort sind die Gewichte über den gesamten Zeitraum gleich.
	
	\begin{figure}[H]
		\centering
		\includegraphics[width=1.2\textwidth]{gewichte_mark.pdf} 
		\caption{Area-Plot der Gewichte für die Markowitz-Optimierung (halbjährlich)}
		\label{fig:area_mark_hj}
	\end{figure}
	
	Die Gewichtsverteilung bei der Markowitz-Optimierung schwankt sehr stark. Außerdem erkennt man, dass wenige Assets hohe Gewichte erhalten. Der Area-Plot legt also nahe dass Markowitz bei jedem Rebalancing große Veränderungen an der Portfoliozusammensetzung vornimmt. Auch wenn in dieser Arbeit die Transaktionskosten nicht berücksichtigt werden, würden die starken Schwankungen bei den Gewichten in der Praxis zu sehr hohen Kosten beim Rebalancing führen.
	
	\begin{figure}[H]
		\centering
		\includegraphics[width=1.2\textwidth]{gewichte_hrp.pdf} 
		\caption{Area-Plot der Gewichte für das HRP-Modell (halbjährlich)}
		\label{fig:area_hrp_hj}
	\end{figure}
	
	Die Gewichtsverteilung beim HRP-Modell ist deutlich gleichmäßiger. Zwar gibt es auch hier Phasen wo die Gewichte stark schwanken, allerdings treten diese starken Schwankungen in Zeiträumen auf wo der Markt aufgrund von Krisen instabil war. Auch fällt sofort auf dass mehr Assets überhaupt mit Gewichten größer null versehen werden. Anders als bei Markowitz scheinen die Gewichte hier eher stabil zu sein, HRP nimmt also weniger starke Änderungen pro Rebalancing vor als Markowitz.
	
	Abschließend werden die Daten zur Performance als Gesamtes betrachtet und für alle drei Rebalancing-Frequenzen. 
	
		\begin{table}[H]
		\centering
		\caption{Vergleich der Performance über verschiedene Rebalancing-Frequenzen}
		\label{tab:performance_vergleich}
		\small
		\begin{tabular}{@{}lccccccccc@{}}
			\toprule
			\textbf{Metrik} & \multicolumn{3}{c}{\textbf{Quartalsweise}} & \multicolumn{3}{c}{\textbf{Halbjährlich}} & \multicolumn{3}{c}{\textbf{Jährlich}} \\ \cmidrule(lr){2-4} \cmidrule(lr){5-7} \cmidrule(lr){8-10}
			& HRP & MinVar & $1/n$ & HRP & MinVar & $1/n$ & HRP & MinVar & $1/n$ \\ \midrule
			Ann. Rendite & 9,61\% & 8,44\% & 10,99\% & 9,63\% & 8,74\% & 10,99\% & 9,32\% & 7,79\% & 10,99\% \\
			Ann. Volatilität & 13,72\% & 12,98\% & 15,26\% & 13,96\% & 13,20\% & 15,26\% & 14,10\% & 13,30\% & 15,26\% \\
			\textbf{Sharpe Ratio} & \textbf{0,701} & 0,650 & 0,720 & \textbf{0,690} & 0,662 & 0,720 & \textbf{0,661} & 0,586 & 0,720 \\
			Max. Drawdown & -47,5\% & -43,7\% & -51,6\% & -47,7\% & -43,3\% & -51,6\% & -49,0\% & -42,8\% & -51,6\% \\
			Avg. Turnover & 0,256 & 0,427 & 0,000 & 0,321 & 0,627 & 0,000 & 0,394 & 0,876 & 0,000 \\ \bottomrule
		\end{tabular}
	\end{table}
	
	Zunächst sieht man dass die Vermutungen über die Rendite, die Volatilität, die Stabilität der Gewichte und den Max. Drawdown aus den Abbildungen \ref{fig:performance}, \ref{fig:area_mark_hj} und \ref{fig:area_hrp_hj} durch die Werte aus der Tabelle bestätigt werden. 
	
	Konkret sieht man das beide Modelle wie erwartet die Volatilität, also das Risiko, verringern und dabei auch die Rendite reduzieren. Das passt auch dazu dass der Max. Drawdown bei beiden Modellen deutlich geringer ist als bei der naiven Diversifikation. Wobei beide Werte von der Markowitz-Optimierung über alle Frequenzen am geringsten sind.
	
	Auch die Stabilität der Gewichte, also der Avg. Turnover ist wie erwartet bei HRP deutlich geringer als bei Markowitz. Für die halbjährliche und die jährliche Rebalancing Frequenz sogar nur ca. halb so groß. Wie bereits erwähnt werden in dieser Arbeit die Transaktionskosten nicht berücksichtigt. Würde man sie allerdings berücksichtigen, dann wäre der Renditeunterschied zwischen HRP und Markowitz noch deutlich größer.
	
	% Hier vielleicht dieses 1/N Puzzle erwähnen, um zu erklären warum das Sharpe-Ration von 1/N so gut ist.
	
	Der vielleicht interessanteste Wert ist das Sharpe-Ratio. Denn es quantifiziert die Effizienz eines Portfolios. Hier sieht man dass die naive Diversifikation über alle Frequenzen das höchste und somit beste Sharpe-Ratio aufweist. Dieses Phänomen konnte Victor DeMiguel bereits in seiner Arbeit von 2009 erklären \parencite[1947]{Puzzle} Auffällig ist auch, dass HRP über alle Frequenzen das bessere Sharpe-Ratio als Markowitz aufweist. Beide Modelle scheinen allerdings von einer höheren Rebalancing-Frequenz zu profitieren, denn je höher die Frequenz ist, desto effizienter werden die HRP- bzw. Markowitz-Portfolios.
	
	% Eingangsfrage beantworten!!!
	
	\textbf{Zusammenfassung:} 
	
	Die Ergebnisse dieser Arbeit bestätigen die in der Theorie bereits geklärten Vorteile des HRP-Algorithmus. Man konnte deutlich stabilere Gewichte bei HRP beobachten. Zudem war das HRP-Portfolio effizienter als das der Markowitz-Optimierung. Zwar hat Markowitz die Volatilität deutlicher gesenkt als HRP, jedoch hat sich dabei die Rendite so sehr reduziert, dass HRP trotzdem das effizientere Porfolio liefert.
	
	Nun zur Eingangsfrage der Arbeit, also zur Frage ob HRP bessere Ergebnisse liefert als die klassische Markowitz-Optimierung. Zunächste kann man festhalten, dass Markowitz tatsächlich Probleme mit der numerischen Stabilität hat. Das Tracking der Konditionszahl hat gezeigt dass diese unter starken Schwankungen leidet. Dieses Problem umgeht HRP komplett. HRP liefert außerdem über alle drei betrachteten Frequenzen ein effizienteres Portfolio als Markowitz. Das kann man durchaus als verbesertes Ergebnis bezeichnen. Es gibt jedoch eine Sache die Markowitz verlässlich besser macht als HRP. Das Risiko wird bei Markowitz mehr reduziert als bei HRP. Sucht man also das Portfolio mit dem geringsten Risiko, dann ist Markowitz die bessere Wahl.
	
	\clearpage
	
	\section{Ausblick}
	
	Diese Untersuchung hat die Vorteile von HRP-Algorithmus gegenüber der Markowitz-Optimierung aufgezeigt. Jedoch ergeben sich aus dieser Arbeit spannende Ansätze für zukünftige Untersuchungen.
	
	\textbf{1. Integration eines realen Kostenmodells:}
	
	HRP lieferte in dieser Untersuchung ein deutlich geringeres Turnover als die Markowitz-Optimierung. Da in dieser Arbeit allerdings sowohl die Transaktionskosten als auch Steuern o.Ä. vernachlässigt wurden, sind die berechneten Werte Bruttowerte. Eine Folgeuntersuchung könnte unter Berücksichtigung dieser Kosten zu dem Ergebnis kommen, dass der Effizienzvorteil von HRP noch größer ist, denn durch das höhere Turnover, würde bei der Markowitz-Optimierung die Nettorendite stärker belastet werden.
	
	\textbf{2. Analyse des 1/N Puzzles:}
	
	Die naive Diversifikation hat interessanterweise über alle betrachteten Frequenzen das effizienteste Portfolio geliefert. In einer Folgestudie könnte man betrachtet ob es eine Portfoliogröße gibt, ab der die Optimierungsverfahren ihren theoretischen Vorteil gegenüber der naiven Diversifikation besser ausspielen können. 
	
	\textbf{3. Variieren weiterer Modell-Parameter:}
	
	In dieser Arbeit wurden drei Frequenzen betrachtet. Allerdings gibt es noch weitere Parameter die man variieren kann. Bspw. wurde hier als Datengrundlage für die Schätzung der Marktparameter eine Zeitraum von einem Jahr gewählt (250 Handelstage). Diesen Zeitraum könnte man in einer späteren Untersuchung variieren, um zu prüfen ob die Modelle von einer größeren Datengrundlage profitieren.
	
	Abschließend lässt sich festhalten, dass HRP eine wichtige Brücke zwischen der modernen Portfoliotheorie und den Anforderungen an numerische Stabilität schlägt. Während die klassische Optimierung ein wertvolles theoretisches Fundament bleibt, bietet der hierarchische Ansatz eine Antwort auf die Herausforderungen verrauschter Finanzmarktdaten.
	
\clearpage
\printbibliography[title={Literaturverzeichnis}]

	
\end{document}
